% GENERATED, edit the French unit or the translation table
\mxunit{automatismes-guide}{User guide (teacher version)}
\subsection*{Principle}
Each Seconde (2G) lesson opens with a sheet of five questions (5 to 8 minutes), aligned with the day's lesson in the calendar:
\begin{enumerate}[label=\textbf{\arabic*.},leftmargin=1.6em]
  \item \textbf{Warm-up}: a Key Stage 3 (cycle~4) prerequisite directly useful for the lesson.
  \item \textbf{Core}: a question on the topic of the day (or of the previous day).
  \item \textbf{Core +}: a second, more demanding question (MCQ, justified true/false, reading code, method).
  \item \textbf{Review}: an out-of-chapter question, drawn from an expectation marked ``to consolidate'' or ``to maintain'' (spaced revision).
  \item \textbf{Thread}: a question targeting an expectation marked ``priority reinforcement'' in the placement test of 29/09/2026.
\end{enumerate}
In the teacher version, each question is preceded by its diagnostic tag (for example \mxtag{EQ}); the expected answer appears in red, followed by a comment or a common pitfall.

\subsection*{Suggested routine (7 min)}
\begin{enumerate}[label=\textbf{\arabic*.},leftmargin=1.6em]
  \item 3 min: individual answers on mini-whiteboards or in the routines notebook, without a calculator unless stated.
  \item 2 min: sharing of answers; have students explain one \emph{wrong} answer and one \emph{correct} answer (``How did you do it?'').
  \item 2 min: targeted correction of questions 4--5 (the diagnostic ones) with the associated method sentence.
\end{enumerate}

\subsection*{Variant for students learning the language of instruction (FLE/FLS)}
\begin{itemize}[leftmargin=1.2em,itemsep=2pt]
  \item \textbf{Instructions}: project the question with the action verb highlighted (\emph{calculate, solve, show, compare}) and a model answer as an example during the first week.
  \item \textbf{Sentence frames}: ``I calculate\ldots{} so\ldots{}'', ``I compare\ldots{} because\ldots{}'', ``The counter-example is\ldots{}''. Display them on a permanent banner.
  \item \textbf{Visual glossary}: keyword sheet (determinant, collinear, median, rate of change, conditional frequency) with a diagram and a translation into the student's home language if useful.
  \item \textbf{Pairs}: 30 seconds of peer rephrasing before answering; the non-native speaker may first answer with a gesture or a sketch, then formalise orally.
  \item \textbf{Written work}: accept symbolic notation before the written sentence, then require the concluding sentence every other time.
\end{itemize}

\subsection*{Error analysis grid}
{\footnotesize\mxarraystretch{1.2}
\begin{tabular}{@{}p{2.2cm} p{4.4cm} p{4.8cm} p{4.4cm}@{}}\toprule
Tags & Observed error & Clue in the answer & Remediation \\ \midrule
EQ, CAL & Sign error when transposing or distributing & $5x-7=2x+8\Rightarrow3x=1$ & Write out each step; check by substituting. \\
DEV, LIT & Forgotten distributivity ($2(x+3)=2x+3$) & Unsimplified result, or wrong for $x=1$ & Numerical test; area of a rectangle. \\
PUI, SUB & Confusing $(-a)^2$ and $-a^2$ & Negative answer for a square & Always use brackets when substituting. \\
GEO, REP & Distance without a square root; $x_A-x_B$ instead of $x_B-x_A$ & $AB=25$ or inconsistent sign & Check the order of magnitude on a sketch. \\
PRO & Adding successive rates; confusing proportion and change & $+\pc{20}$ then $-\pc{10}=+\pc{10}$ & Use multipliers. \\
PRB, LEC & Confusing $P_A(B)$ and $P_B(A)$ & Wrong reference total & Circle the total of the subset after ``given that''. \\
REG, IMG & Confusing image and preimage; reading reversed & Answer given on the wrong axis & Colour-code: abscissa = input, ordinate = output. \\
DEM, LIT & An example taken as a proof & ``It works for $n=3$'' & Distinguish testing, counter-example and proof. \\
ALG & Running code without tracking the variables & Incorrect final value & Trace table (variable / value / step). \\ \bottomrule
\end{tabular}}

\subsection*{Marking scheme}
The routines are formative: no marking scheme is provided. A detailed scheme by competence (Investigate, Model, Represent, Reason, Calculate, Communicate) can be added if one of these sheets is used for assessment.

\subsection*{Sources and points to watch}
\begin{itemize}[leftmargin=1.2em,itemsep=2pt]
  \item Mathematics programme for general and technological Seconde, order of 26-2-2026 (NOR MENE2602914A), BO no.~14 of 2 April 2026: section ``Automatismes'' and the parts ``Numbers and calculation'', ``Geometry'', ``Functions'', ``Statistics and probability'', ``Set vocabulary and logic''.
  \item Calendar: consulted on 06/10/2026; the order of topics follows the calendar and not the dates in the progression document, which differ (e.g.\ functions: 3--15 Dec.\ 2026 in the calendar).
  \item Exercises: \textbf{created for this document} with your permission; all numerical and algebraic values were recomputed by script (see \texttt{tools/}), and the reasoning was proofread. The data in the ``rainfall'' table (sheet~35) are fictitious.
  \item Quartile convention: smallest value such that at least $\pc{25}$ (resp.\ $\pc{75}$) of the data are less than or equal to it (rank $\lceil n/4\rceil$, $\lceil3n/4\rceil$) --- \textbf{[VERIFY]} coordinate with the conventions of the textbook used.
  \item Questions on solids, nets and perspective (tags SOL, PER) belong to Key Stage 3 (cycle~4): the Seconde programme has no dedicated chapter; they serve the diagnostic.
  \item Sheet~77 (question~3): notion of rate of change in anticipation of Première, outside the Seconde chapters.
  \item List of automatismes assessable in the early examination (BO no.~24 of 12 June 2025, 2027 session): \textbf{[VERIFY --- check Eduscol]} for the 2028 session.
\end{itemize}
